\subsection{The Dependency Map} \label{subsec:the-dependency-map}
\par In a symbolic expression we sometimes see variables defined in the context of other variables.
This primarily happens when a variable is declared inside a Kleene closure. 
For example the expression $a \asgnarrow (b \asgnarrow b)*_{a}$ will match any list of elements. 
The elements won't necessarily all be equal even though they were all declared using the same variable $b$.
This is why we use the context so that $b$ can take on a different value for each iteration of $a$.
However, if $b$ is referenced without $a$ being in context, it becomes ambiguous which $b$ should be used.
This is why we introduce the dependency map $\deps$. 
The dependency map is a dictionary that keeps track of which variables depend on which other variables. 
In the example above, $a$ doesn't depend on anything, and $b$ depends on $a$. 
This means that it is illegal to reference $b$ unless $a$ is in the current context.

\subsubsection{Dependency Map Computation} \label{subsubsec:dependency-map-computation}
The dependency map is computer for a group of symbolic expressions that represent a function. 
The computation traverses the first expression keeping track of the context as it goes. 
When it comes across an assignment for some variable $\lab$, it adds an entry like so $\deps[\pathvar \rightarrow \lab]$.
Once it finishes traversing the first tree, it has a dependency map for the first tree but not the whole function. So it
threads the result from the first traversal and repeats the process with the subsequent expressions. A variable may be used
in multiple contexts throughout all the expressions, but the first use of the variable is the one that ends up in the final
dependency map.

\subsubsection{Dependency Map Use} \label{subsubsec:dependency-map-use}
The dependency map is used for environment lookups and updates. 
First, recall that every key in the environment is a context-label pair that gets mapped to a value.
So whenever a lookup or update occurs, the first thing that happens is we verify that all the label's
depencies are satisfied by the context that was passed in. Then we filter out all variables from the
context that are not in the label's dependencies. This gives us $\pathvar'$ which contains only the
relevant context for the label. This relevant context is what is actually used to lookup values and
insert pairs into the environment. 
\evan{Should I give an example here to demonstrate why these steps are necessary?}
